\section{Protocol Design}
\label{sec:protocol}

We present the Lux consensus protocol family as a progression of increasingly robust mechanisms, each building on the prior.

\subsection{System Model}

Consider a network of $n$ validators $\{v_1, \ldots, v_n\}$ with stake weights $\{s_1, \ldots, s_n\}$ where $\sum s_i = S$. An adversary controls validators with combined stake $f < S/4$. Communication is partially synchronous: messages between correct validators arrive within bound $\Delta$ after GST (Global Stabilization Time).

\subsection{Slush: Non-BFT Metastable Protocol}

Slush is the simplest protocol in the family. Each validator maintains a color (preference) $c \in \{R, B, \bot\}$.

\begin{algorithm}[h]
\caption{Slush Protocol}
\begin{algorithmic}[1]
\State Initialize $c \gets \bot$
\Upon{receiving transaction $\mathit{tx}$ with color $c'$}
    \If{$c = \bot$} $c \gets c'$ \EndIf
    \For{round $r = 1$ to $m$}
        \State Sample $k$ validators uniformly at random
        \State Query sampled validators for their color
        \State Let $c_{\max}$ be the majority color among responses
        \If{$|\{j : c_j = c_{\max}\}| \geq \alpha k$} $c \gets c_{\max}$ \EndIf
    \EndFor
\EndUpon
\end{algorithmic}
\end{algorithm}

The parameter $\alpha \in (0.5, 1]$ is the \emph{confidence threshold}. Higher $\alpha$ increases resistance to Byzantine nodes but requires more rounds for convergence.

\subsection{Snowflake: BFT with Conviction Counter}

Snowflake augments Slush with a \emph{conviction counter} $\mathit{cnt}$ that tracks consecutive rounds agreeing with the current preference. A decision is made when $\mathit{cnt}$ reaches threshold $\beta$.

\begin{algorithm}[h]
\caption{Snowflake Protocol}
\begin{algorithmic}[1]
\State Initialize $c \gets \bot$, $\mathit{cnt} \gets 0$
\Upon{query round}
    \State Sample $k$ validators, collect responses
    \State Let $c_{\max}$ be the $\alpha$-majority color
    \If{$c_{\max} = c$}
        \State $\mathit{cnt} \gets \mathit{cnt} + 1$
    \Else
        \State $c \gets c_{\max}$, $\mathit{cnt} \gets 1$
    \EndIf
    \If{$\mathit{cnt} \geq \beta$} \textbf{decide} $c$ \EndIf
\EndUpon
\end{algorithmic}
\end{algorithm}

The conviction counter provides Byzantine fault tolerance: an adversary must sustain a majority response for $\beta$ consecutive rounds to alter a validator's decision.

\subsection{Snowball: Confidence with Historical Memory}

Snowball extends Snowflake by maintaining per-color \emph{confidence counters} $d[c]$, providing historical memory:

\begin{algorithm}[h]
\caption{Snowball Protocol}
\begin{algorithmic}[1]
\State Initialize $c \gets \bot$, $d[\cdot] \gets 0$, $\mathit{cnt} \gets 0$
\Upon{query round}
    \State Sample $k$ validators, collect responses
    \State Let $c_{\max}$ be the $\alpha$-majority color
    \State $d[c_{\max}] \gets d[c_{\max}] + 1$
    \If{$d[c_{\max}] > d[c]$} $c \gets c_{\max}$ \EndIf
    \If{consecutive for $c$} $\mathit{cnt} \gets \mathit{cnt} + 1$ \Else{} $\mathit{cnt} \gets 1$ \EndIf
    \If{$\mathit{cnt} \geq \beta$} \textbf{decide} $c$ \EndIf
\EndUpon
\end{algorithmic}
\end{algorithm}

The per-color counters prevent an adversary from resetting progress by briefly flipping the majority. A preference changes only when cumulative evidence exceeds the current choice.

\subsection{Avalanche: DAG-Based Multi-Transaction Consensus}

Avalanche extends Snowball to operate over a DAG of transactions, amortizing consensus across the structure:

\begin{itemize}
    \item Each transaction $T$ is a vertex in the DAG with parent edges encoding dependencies.
    \item A \emph{conflict set} $\mathcal{C}_T$ contains all transactions conflicting with $T$ (e.g., double-spends).
    \item Each vertex maintains a Snowball instance. Querying a vertex transitively supports all its ancestors.
    \item A transaction is \emph{accepted} when its Snowball instance decides and all ancestors are accepted.
\end{itemize}

This DAG structure enables:
\begin{enumerate}
    \item \textbf{Throughput amortization}: One query round implicitly votes on all ancestor transactions.
    \item \textbf{Efficient conflict resolution}: Only conflicting transactions require explicit consensus; compatible transactions are accepted transitively.
    \item \textbf{Pipelining}: New transactions can be proposed before prior ones finalize, enabling continuous throughput.
\end{enumerate}

\subsection{Stake-Weighted Sampling}

In the permissionless setting, Sybil resistance requires weighting validator influence by stake. We replace uniform random sampling with stake-weighted sampling:

\begin{equation}
    P(\text{sample } v_i) = \frac{s_i}{\sum_{j=1}^{n} s_j}
\end{equation}

The $\alpha$-majority threshold is correspondingly defined over stake:
\begin{equation}
    \alpha\text{-majority for } c \iff \sum_{j: c_j = c} s_j \geq \alpha \sum_{j=1}^{k} s_j
\end{equation}

\subsection{Multi-Chain Architecture}

Lux employs the consensus protocol family across multiple chains:

\begin{description}
    \item[X-Chain] Uses the Avalanche DAG protocol for UTXO-based asset transfers. Optimized for throughput via DAG amortization.
    \item[C-Chain] Uses Snowball with linear ordering for EVM-compatible smart contracts. Supports Solidity, deterministic execution.
    \item[P-Chain] Uses Snowball for platform operations: validator management, staking, subnet creation, and governance.
    \item[Subnets] Independent validator subsets running custom VMs with shared or independent consensus parameters.
\end{description}

All chains share a common validator set (with subnet specialization) and a unified staking mechanism, while maintaining independent state machines and execution environments.
